% theme: stimulus_and_response
\MajorQuestion{刺激と反応}
\SetRowSpaceQanda

\Instruction{次の問いに答えなさい。}
\QQ{外界からの刺激を受けとる器官を\NamedBlank{sensoryorgan}という。}{感覚器官}
\QQ{筋肉の両端にあり、筋肉を骨につなぐ部分を何というか。}{けん}
\QQ{脳や脊髄をまとめて何というか。}{中枢神経}
\QQ{目で、光の刺激を受けとる細胞がある部分を何というか。}{網膜}
\QQ{手やあしなど、動物の体を動かす器官を何というか。}{運動器官}

\Instruction{\strut}
\QQ{\RefBlank{sensoryorgan}からの刺激の信号を、脳や脊髄に伝える神経を何というか。}{感覚神経}
\QQ{骨と骨のつなぎ目を何というか。}{関節}
\QQ{熱いものにふれて、思わず手を引っこめるような反応を何というか。}{反射}
\QQ{耳で、音によって振動する膜を何というか。}{鼓膜}
\QQ{体を支えたり、内臓や脳などを保護したりする、多数の骨の組み合わせを\NamedBlank{skeleton}という。}{骨格}

\Instruction{\strut}
\QQ{脳や脊髄から枝分かれして全身に広がっている神経を何というか。}{末しょう神経}
\QQ{皮膚で受けとる刺激を、3つ答えなさい。}{接触、圧力、あたたかさ、冷たさ、痛みのうち3つ}
\QQ{意識して起こる反応では、刺激の信号は\RefBlank{sensoryorgan}から脳に伝わった後、どこを通って運動器官へ伝わるか。}{脊髄}
\QQ{鼻には何の刺激を受けとる細胞があるか。}{におい}
\QQ{刺激や命令の信号を伝える神経や、刺激を判断して命令を出す脳や脊髄をまとめて何というか。}{神経系}

\Instruction{\strut}
\QQ{うでを曲げるとき、内側と外側の筋肉はそれぞれどうなるか。}{内側が収縮し、外側がゆるむ}
\QQ{反射では、刺激の信号は脳を経由するか。}{経由しない}
\QQ{脳や脊髄からの命令の信号を、運動器官に伝える神経を何というか。}{運動神経}
\QQ{反射では、刺激の信号は\RefBlank{sensoryorgan}から何を経由して運動器官に伝わるか。}{脊髄}
\QQ{反射は、刺激を受けてから反応するまでの時間が短い。このことは何に役立つか。}{危険から体を守ること}
